% appendices/appendix-b-tugas-portofolio.tex
\chapter{Contoh Tugas, Portofolio, dan Lembar Kerja}
\label{app:lembar-kerja}

Lampiran ini memuat bentuk-bentuk yang dipakai sepanjang perkuliahan dalam
keadaan siap dicetak. Berkas aslinya tersedia pada repositori mata kuliah
dalam bentuk yang dapat disunting; yang disajikan di sini adalah bentuk dan
isinya, agar Anda mengetahui apa yang diharapkan sebelum mengerjakannya.

\section{Format Laporan Praktikum}

Setiap praktikum pada Bab~III sampai Bab~XV menghasilkan satu berkas laporan.
Laporan itu tidak perlu panjang; satu halaman sudah cukup, asalkan memuat
keempat bagian pada Tabel~\ref{tab:format-laporan}.

\begin{table}[htbp]
\centering
\caption{Format laporan praktikum satu halaman}
\label{tab:format-laporan}
{\setlength{\tabcolsep}{4pt}\footnotesize%
\begin{tabularx}{\textwidth}{@{}L{3.0cm} Y@{}}
\toprule
\headrow \textbf{Bagian} & \textbf{Isi, dan mengapa diperlukan} \\
\midrule
Tujuan & Satu atau dua kalimat yang menyebut Sub-CPMK dan hasil yang hendak
dicapai. Bagian ini memudahkan Anda memeriksa sendiri apakah praktikumnya
berhasil. \\
\addlinespace
Yang dikerjakan & Uraian singkat langkah yang diambil, termasuk perubahan yang
tidak ada pada petunjuk. Bagian ini yang membedakan laporan hasil kerja dari
salinan petunjuk. \\
\addlinespace
Keluaran & Tangkapan layar atau cuplikan keluaran program, beserta angkanya
apabila ada. Keluaran yang disertai angka dapat diperiksa ulang. \\
\addlinespace
Catatan & Masalah yang ditemui, dugaan penyebabnya, dan cara mengatasinya.
Bagian ini yang paling berguna bagi Anda pada praktikum berikutnya. \\
\bottomrule
\end{tabularx}}
\end{table}

\section{Format Portofolio}

Portofolio adalah kumpulan seluruh artefak praktikum pada satu repositori.
Susunan yang dianjurkan diperlihatkan pada Tabel~\ref{tab:format-portofolio}.
Susunan ini meniru susunan sumber buku ini, sehingga mahasiswa yang mengikuti
seluruh praktikum akan memperoleh kebiasaan yang sama.

\begin{table}[htbp]
\centering
\caption{Susunan direktori portofolio}
\label{tab:format-portofolio}
{\setlength{\tabcolsep}{4pt}\footnotesize%
\begin{tabularx}{\textwidth}{@{}L{4.4cm} Y@{}}
\toprule
\headrow \textbf{Direktori} & \textbf{Isi} \\
\midrule
\code{code/python/} & Seluruh berkas Python dari Bab~III sampai Bab~XI, dengan
nama yang menyebut babnya. \\
\addlinespace
\code{code/unity/} & Seluruh berkas C\# dari Bab~XII sampai Bab~XVI. \\
\addlinespace
\code{laporan/} & Laporan praktikum satu halaman untuk setiap pertemuan. \\
\addlinespace
\code{aset/} & Aset yang dipakai, beserta berkas \code{SUMBER.md} yang memuat
sumber dan lisensi setiap aset. \\
\addlinespace
\code{proyek/} & Proyek akhir: \istilah{build}, kode sumber, laporan
pengembangan, lembar \istilah{playtest}, dan \istilah{log bug}. \\
\bottomrule
\end{tabularx}}
\end{table}

Aset berukuran besar sebaiknya tidak disimpan langsung di dalam repositori.
Simpanlah pada tempat penyimpanan berkas, lalu catat alamatnya pada berkas
\code{SUMBER.md}.

\section{Lembar \istilah{Playtest} Siap Cetak}

Lembar pada Tabel~\ref{tab:lembar-playtest-cetak} dipakai pada sesi
\istilah{playtest} antar kelompok. Isilah satu lembar untuk setiap penguji.

\begin{table}[htbp]
\centering
\caption{Lembar \istilah{playtest} siap cetak}
\label{tab:lembar-playtest-cetak}
{\setlength{\tabcolsep}{4pt}\footnotesize%
\begin{tabularx}{\textwidth}{@{}L{4.0cm} Y@{}}
\toprule
\headrow \textbf{Butir} & \textbf{Isian} \\
\midrule
Nama penguji dan kelompok & \\[1.3em]
\midrule
Pengalaman bermain sejenis & \\[1.3em]
\midrule
Nama \istilah{build} dan tanggalnya & \\[1.3em]
\midrule
Waktu mulai dan waktu selesai & \\[1.3em]
\midrule
Tujuan sesi (bagian yang diragukan) & \\[1.3em]
\midrule
Titik berhenti lebih dari 10 detik & \\[2.0em]
\midrule
Keadaan yang tidak seharusnya terjadi & \\[2.0em]
\midrule
Pendapat penguji, dicatat apa adanya & \\[2.0em]
\midrule
Catatan pengamat sendiri & \\[2.0em]
\bottomrule
\end{tabularx}}
\end{table}

\section{Log \istilah{Bug} Siap Cetak}

Log pada Tabel~\ref{tab:log-bug-cetak} memuat kolom yang sama dengan
Tabel~\ref{tab:log-bug} pada Bab~XVI, dalam bentuk yang dapat diisi langsung.

\begin{table}[htbp]
\centering
\caption{Log \istilah{bug} siap cetak}
\label{tab:log-bug-cetak}
{\setlength{\tabcolsep}{4pt}\footnotesize%
\begin{tabularx}{\textwidth}{@{}C{0.7cm} C{1.6cm} Y C{1.3cm} C{1.4cm}@{}}
\toprule
\headrow \textbf{No} & \textbf{Jenis} & \textbf{Urutan yang menyebabkannya, dan hasil yang salah} & \textbf{Berat} & \textbf{Status} \\
\midrule
1 & & & & \\[2.2em]
\midrule
2 & & & & \\[2.2em]
\midrule
3 & & & & \\[2.2em]
\midrule
4 & & & & \\[2.2em]
\midrule
5 & & & & \\[2.2em]
\bottomrule
\end{tabularx}}
\end{table}

Jenis temuan yang dipakai adalah mekanik, integrasi, tampilan, dan
\istilah{soft-lock}. Beratnya adalah berat, sedang, atau ringan, menurut
Tabel~\ref{tab:berat-temuan}. Statusnya adalah belum, sedang dikerjakan, atau
selesai beserta tanggalnya.

\section{Daftar Periksa Penyerahan Proyek Akhir}

Serahkan seluruh butir pada Tabel~\ref{tab:penyerahan} dalam satu
\istilah{build} dan satu repositori. Berkas yang hilang pada hari penyerahan
tidak dapat disusulkan tanpa persetujuan dosen pengampu.

\begin{table}[htbp]
\centering
\caption{Daftar periksa penyerahan Proyek Akhir}
\label{tab:penyerahan}
{\setlength{\tabcolsep}{4pt}\footnotesize%
\begin{tabularx}{\textwidth}{@{}C{0.8cm} L{6.2cm} Y@{}}
\toprule
\headrow \textbf{No} & \textbf{Kelengkapan} & \textbf{Keterangan} \\
\midrule
1 & \istilah{Build} yang dapat dijalankan & Diuji pada komputer lain, bukan
pada komputer pembuatnya. \\
\addlinespace
2 & Kode sumber \istilah{project} & Seluruh berkas yang diperlukan untuk
membangun ulang. \\
\addlinespace
3 & Ringkasan desain (GDD) & Beserta catatan perubahan yang disengaja. \\
\addlinespace
4 & Laporan pengembangan & Memuat keadaan sebelum dan sesudah setiap
perbaikan. \\
\addlinespace
5 & Lembar \istilah{playtest} terisi & Sekurang-kurangnya dua penguji dari
kelompok lain. \\
\addlinespace
6 & \istilah{Log bug} & Beserta status penanganan setiap temuan. \\
\addlinespace
7 & Daftar aset dan lisensinya & Memuat sumber setiap aset yang bukan karya
sendiri. \\
\addlinespace
8 & Lembar refleksi akhir semester & Menurut Bagian~\ref{sec:refleksi-akhir}. \\
\bottomrule
\end{tabularx}}
\end{table}

\section{Contoh Tugas yang Dapat Diberikan Dosen}

Tabel~\ref{tab:contoh-tugas} memuat contoh tugas tambahan yang dapat dipakai
apabila suatu CPMK belum mencapai ambang batas. Setiap tugas menunjuk pada
bagian buku yang menjadi sumbernya, sehingga mahasiswa dapat mengerjakannya
tanpa bahan tambahan.

\begin{table}[htbp]
\centering
\caption{Contoh tugas tambahan dan bagian buku yang menjadi sumbernya}
\label{tab:contoh-tugas}
{\setlength{\tabcolsep}{4pt}\footnotesize%
\begin{tabularx}{\textwidth}{@{}L{2.0cm} L{4.6cm} Y@{}}
\toprule
\headrow \textbf{CPMK} & \textbf{Bentuk tugas} & \textbf{Sumber} \\
\midrule
\CPMK{1} & Mengukur laju bingkai pada tiga nilai pembatas, lalu menjelaskan
hubungannya dengan beban kerja & Bab~IV, dan \code{ch04_ukur_laju.py} \\
\addlinespace
\CPMK{1} & Membandingkan gerak berbasis jumlah bingkai dengan gerak berbasis
\istilah{delta time} pada dua ukuran jendela & Bab~IV dan Bab~V \\
\addlinespace
\CPMK{2} & Menelusuri urutan keadaan NPC pada satu skenario, dan menunjukkan
transisi yang tidak sah & Bab~IX, dan \code{ch09_fsm.py} \\
\addlinespace
\CPMK{2} & Membandingkan banyak petak yang diperiksa oleh dua heuristik pada
peta yang sama & Bab~X, dan \code{ch10_astar.py} \\
\addlinespace
\CPMK{3} & Menjelaskan gejala fisika pada sebuah \istilah{scene}, dan
menentukan komponen penyebabnya & Bab~XIII \\
\addlinespace
\CPMK{4} & Menyusun daftar peristiwa permainan beserta bunyi yang seharusnya
memicunya & Bab~XV \\
\addlinespace
\CPMK{4} & Menjalankan \code{ch16_uji_mandiri.py} atas level sendiri, lalu
memperbaiki tempat yang dilaporkan & Bab~XVI \\
\bottomrule
\end{tabularx}}
\end{table}
